\chapter{Function Approximation and the Deadly Triad}

Chapters 3--6 stored value in a \textbf{table}, one entry $V(s)$ or $Q(s,a)$ per state or pair, updated only at visited keys. Chapter 6 then specified how to choose $A_t$ from that table ($\varepsilon$-greedy or softmax). Tabular control is fine when $\mathcal{S}$ and $\mathcal{A}$ are small and discrete, as in Chapters 3--6. The same $s$ returns, and the same cell is updated. Many tasks have continuous states, images, or huge discrete spaces. Then a separate number for every $(s,a)$ runs into two problems: the table does not fit in memory, and a later observed state is almost never that same $s$, so the stored number is not used again.

\textbf{Function approximation} addresses both. In the tabular case the object we store is a whole map---either $s\mapsto V(s)$ or $(s,a)\mapsto Q(s,a)$---one free number per key. Replace that map by a function of a weight vector $\mathbf{w}\in\mathbb{R}^d$ with components $w_1,\ldots,w_d$:

\[
\hat{V}(s;\mathbf{w})\approx V(s)
\qquad\text{or}\qquad
\hat{Q}(s,a;\mathbf{w})\approx Q(s,a),
\]

with $d$ much smaller than the number of keys (or finite $d$ even when $\mathcal{S}$ is continuous). Only the $d$ weights are stored and updated. For action-values we use a \textbf{feature map} $\mathbf{x}$: given $(s,a)$, it returns a feature vector $\mathbf{x}(s,a)\in\mathbb{R}^d$ of the same length as $\mathbf{w}$. The map itself is fixed---the same pair always produces the same vector; different pairs may produce different vectors. The linear case in (7.1) is the inner product $\hat{Q}(s,a;\mathbf{w})=\mathbf{w}^\top\mathbf{x}(s,a)$. (State-values use an analogous $\mathbf{x}(s)\in\mathbb{R}^d$ and $\hat{V}(s;\mathbf{w})=\mathbf{w}^\top\mathbf{x}(s)$.) At time $t$ one writes $\mathbf{x}_t=\mathbf{x}(S_t,A_t)$ for that step's feature vector. Changing $\mathbf{w}$ changes the prediction at every input whose features use those weights, so experience at one key transfers to others---that is generalization. The rest of this section and Chapter 8 develop the $Q$ form used for control; the $V$ form is the same idea for prediction. Exploration is unchanged: $\varepsilon$-greedy or softmax still choose $A_t$ from predicted action-values (Chapter 6). What changes is the representation and the update.

This section develops that idea in two steps. First, Section 7.1 lifts tabular control to a linear approximator: off-policy Q-learning FA (7.3) and on-policy SARSA FA (7.4). Function approximation itself is neither on-policy nor off-policy---that label belongs to the update. Second, Section 7.2 discusses the \textbf{deadly triad}, in which off-policy semi-gradient FA can diverge, while on-policy linear FA is known to be stable. Chapter 8 then introduces the practical fixes used in deep Q-learning.

\section{Linear Function Approximation}

As discussed at the start of this section, a separate table entry for every $(s,a)$ is often infeasible, so we represent $Q$ by a function of weights. The linear form is

\begin{equation}
\hat{Q}(s,a;\mathbf{w})
=
\mathbf{w}^\top \mathbf{x}(s,a)
=
\sum_{i=1}^{d} w_i\, x_i(s,a),
\tag{7.1}
\end{equation}

with a fixed feature map $\mathbf{x}(s,a)\in\mathbb{R}^d$. Only $\mathbf{w}$ is learned. Differentiating (7.1) gives

\begin{equation}
\nabla_{\mathbf{w}} \hat{Q}(s,a;\mathbf{w})
=
\mathbf{x}(s,a).
\tag{7.2}
\end{equation}

For state-values the same construction uses features of the state alone, $\hat{V}(s;\mathbf{w})=\mathbf{w}^\top\mathbf{x}(s)$.

We still learn online over episodes. At each time step $t$ of an episode, the agent is at $(S_t,A_t)$, observes a transition, and updates $\mathbf{w}$ once. Write $\mathbf{x}_t=\mathbf{x}(S_t,A_t)$ and $\hat{Q}_t=\hat{Q}(S_t,A_t;\mathbf{w})=\mathbf{w}^\top\mathbf{x}_t$. The TD control updates from Chapter 6 carry over with the table $Q$ replaced by the approximator $\hat{Q}$. \textbf{On-policy} (SARSA) is

\[
\begin{aligned}
\hat{Q}(S_t,A_t;\mathbf{w})
&\leftarrow
\hat{Q}(S_t,A_t;\mathbf{w})
+
\alpha\bigl[R_{t+1}+\gamma\hat{Q}(S_{t+1},A_{t+1};\mathbf{w})-\hat{Q}(S_t,A_t;\mathbf{w})\bigr]
\\
&=
\hat{Q}(S_t,A_t;\mathbf{w})
+
\alpha\bigl(U_t-\hat{Q}(S_t,A_t;\mathbf{w})\bigr),
\end{aligned}
\]

with $U_t=R_{t+1}+\gamma\hat{Q}(S_{t+1},A_{t+1};\mathbf{w})$. \textbf{Off-policy} (Q-learning) is

\[
\begin{aligned}
\hat{Q}(S_t,A_t;\mathbf{w})
&\leftarrow
\hat{Q}(S_t,A_t;\mathbf{w})
+
\alpha\bigl[R_{t+1}+\gamma\max_{a'}\hat{Q}(S_{t+1},a';\mathbf{w})-\hat{Q}(S_t,A_t;\mathbf{w})\bigr]
\\
&=
\hat{Q}(S_t,A_t;\mathbf{w})
+
\alpha\bigl(U_t-\hat{Q}(S_t,A_t;\mathbf{w})\bigr),
\end{aligned}
\]

with $U_t=R_{t+1}+\gamma\max_{a'}\hat{Q}(S_{t+1},a';\mathbf{w})$. (If $S_{t+1}$ is terminal, take the $\gamma$ term as $0$ in either case.) The two methods differ only in that backup $U_t$.

Those displays state the intended change in $\hat{Q}$ at the visited pair. There is no separate table entry to write into: $\hat{Q}$ is determined by $\mathbf{w}$, so the update is realized by changing $\mathbf{w}$. The squared gap to close is the one-step objective

\[
J_t(\mathbf{w})
=
\frac{1}{2}\bigl(U_t-\hat{Q}(S_t,A_t;\mathbf{w})\bigr)^2
=
\frac{1}{2}\bigl(U_t-\mathbf{w}^\top\mathbf{x}_t\bigr)^2.
\]

Because $U_t$ bootstraps $\hat{Q}(\cdot;\mathbf{w})$, it depends on $\mathbf{w}$. The gradient of $J_t$ is

\[
\nabla_{\mathbf{w}} J_t(\mathbf{w})
=
\bigl(U_t-\mathbf{w}^\top\mathbf{x}_t\bigr)
\bigl(\nabla_{\mathbf{w}} U_t-\mathbf{x}_t\bigr).
\]

\textbf{Full-gradient} descent uses that expression as written. \textbf{Semi-gradient} sets $\nabla_{\mathbf{w}} U_t=\mathbf{0}$ when differentiating, then $\mathbf{w}\leftarrow\mathbf{w}+\alpha\bigl(U_t-\mathbf{w}^\top\mathbf{x}_t\bigr)\,\mathbf{x}_t$. How $U_t$ is formed decides on-policy vs off-policy. The next two subsections spell each case out.

\subsection{Q-learning}

Q-learning is \textbf{off-policy} because the $\max$ backup learns about a greedy continuation while behavior may explore (Chapter 6). Take the off-policy $U_t$ from above,

\[
U_t=R_{t+1}+\gamma\max_{a'}\hat{Q}(S_{t+1},a';\mathbf{w}).
\]

For the greedy action $a^*=\arg\max_{a'}\hat{Q}(S_{t+1},a';\mathbf{w})$ (break ties arbitrarily), $U_t=R_{t+1}+\gamma\mathbf{w}^\top\mathbf{x}(S_{t+1},a^*)$ when $S_{t+1}$ is nonterminal. Differentiating and using (7.2) gives $\nabla_{\mathbf{w}} U_t=\gamma\mathbf{x}(S_{t+1},a^*)$ (since $R_{t+1}$ does not depend on $\mathbf{w}$). If $S_{t+1}$ is terminal, $U_t=R_{t+1}$ and $\nabla_{\mathbf{w}} U_t=\mathbf{0}$.

Full-gradient descent on $J_t$ is

\[
\begin{aligned}
\mathbf{w}
&\leftarrow
\mathbf{w}
-
\alpha\,\nabla_{\mathbf{w}} J_t(\mathbf{w})
\\
&=
\mathbf{w}
-
\alpha\,
\bigl(U_t-\mathbf{w}^\top\mathbf{x}_t\bigr)
\bigl(\nabla_{\mathbf{w}} U_t-\mathbf{x}_t\bigr).
\end{aligned}
\]

Semi-gradient sets $\nabla_{\mathbf{w}} U_t=\mathbf{0}$ and yields

\begin{equation}
\mathbf{w}
\leftarrow
\mathbf{w}
+
\alpha
\Bigl[
R_{t+1}
+
\gamma\max_{a'}\bigl(\mathbf{w}^\top\mathbf{x}(S_{t+1},a')\bigr)
-
\mathbf{w}^\top\mathbf{x}_t
\Bigr]
\mathbf{x}_t.
\tag{7.3}
\end{equation}

Equivalently, $\mathbf{w}\leftarrow\mathbf{w}+\alpha\bigl[R_{t+1}+\gamma\max_{a'}\hat{Q}(S_{t+1},a';\mathbf{w})-\hat{Q}(S_t,A_t;\mathbf{w})\bigr]\mathbf{x}(S_t,A_t)$. This keeps the $\max$ backup, so it is off-policy.

Semi-gradient ignores the dependence of $U_t$ on $\mathbf{w}$ when differentiating $J_t$. That does not freeze the target across learning: after $\mathbf{w}$ is updated, the next transition builds a new $U_t$ from the new $\hat{Q}$. So the bootstrap still \textbf{moves} every step.

\begin{implnote}[Algorithm (full-gradient Q-learning with linear FA)]
Input: feature map $\mathbf{x}(s,a)$; step size $\alpha$; discount $\gamma$; exploration (e.g. $\varepsilon$-greedy).

Output: $\mathbf{w}$ such that $\hat{Q}(s,a;\mathbf{w})=\mathbf{w}^\top\mathbf{x}(s,a)\approx Q^*(s,a)$; greedy policy $\pi(s)=\arg\max_a\hat{Q}(s,a;\mathbf{w})$.

\begin{enumerate}
\item Initialize $\mathbf{w}$ arbitrarily (e.g. $\mathbf{0}$).
\item Loop for each episode:
\begin{itemize}
\item Initialize $S$.
\item For each step of the episode until $S$ is terminal:
\begin{itemize}
\item Choose $A$ from $S$ using a policy based on $\hat{Q}(\cdot,\cdot;\mathbf{w})$ (e.g. $\varepsilon$-greedy).
\item Take $A$; observe $R$, $S'$.
\item $\mathbf{x}\leftarrow\mathbf{x}(S,A)$.
\item If $S'$ is terminal:
\begin{itemize}
\item $U\leftarrow R$
\item $\mathbf{g}_U\leftarrow\mathbf{0}$
\end{itemize}
\item Else:
\begin{itemize}
\item $a^*\leftarrow\arg\max_{a'}\mathbf{w}^\top\mathbf{x}(S',a')$
\item $U\leftarrow R+\gamma\mathbf{w}^\top\mathbf{x}(S',a^*)$
\item $\mathbf{g}_U\leftarrow\gamma\mathbf{x}(S',a^*)$
\end{itemize}
\item $\mathbf{w}\leftarrow\mathbf{w}-\alpha\bigl(U-\mathbf{w}^\top\mathbf{x}\bigr)(\mathbf{g}_U-\mathbf{x})$.
\item $S\leftarrow S'$.
\end{itemize}
\end{itemize}
\item Return $\mathbf{w}$ and $\pi(s)=\arg\max_a\hat{Q}(s,a;\mathbf{w})$.
\end{enumerate}

\end{implnote}

\begin{implnote}[Algorithm (semi-gradient Q-learning with linear FA)]
Input: feature map $\mathbf{x}(s,a)$; step size $\alpha$; discount $\gamma$; exploration (e.g. $\varepsilon$-greedy).

Output: $\mathbf{w}$ such that $\hat{Q}(s,a;\mathbf{w})=\mathbf{w}^\top\mathbf{x}(s,a)\approx Q^*(s,a)$; greedy policy $\pi(s)=\arg\max_a\hat{Q}(s,a;\mathbf{w})$.

\begin{enumerate}
\item Initialize $\mathbf{w}$ arbitrarily (e.g. $\mathbf{0}$).
\item Loop for each episode:
\begin{itemize}
\item Initialize $S$.
\item For each step of the episode until $S$ is terminal:
\begin{itemize}
\item Choose $A$ from $S$ using a policy based on $\hat{Q}(\cdot,\cdot;\mathbf{w})$ (e.g. $\varepsilon$-greedy).
\item Take $A$; observe $R$, $S'$.
\item $\mathbf{x}\leftarrow\mathbf{x}(S,A)$.
\item If $S'$ is terminal:
\begin{itemize}
\item $U\leftarrow R$
\end{itemize}
\item Else:
\begin{itemize}
\item $U\leftarrow R+\gamma\max_{a'}\mathbf{w}^\top\mathbf{x}(S',a')$
\end{itemize}
\item $\mathbf{w}\leftarrow\mathbf{w}+\alpha\bigl(U-\mathbf{w}^\top\mathbf{x}\bigr)\,\mathbf{x}$.
\item $S\leftarrow S'$.
\end{itemize}
\end{itemize}
\item Return $\mathbf{w}$ and $\pi(s)=\arg\max_a\hat{Q}(s,a;\mathbf{w})$.
\end{enumerate}

\end{implnote}

\subsection{SARSA}

SARSA is \textbf{on-policy} because the backup uses the next action $A_{t+1}$ from the same behavior that chose $A_t$ (Chapter 6). Take the on-policy $U_t$ from above,

\[
U_t
=
R_{t+1}+\gamma\hat{Q}(S_{t+1},A_{t+1};\mathbf{w})
=
R_{t+1}+\gamma\mathbf{w}^\top\mathbf{x}(S_{t+1},A_{t+1})
\]

(with the $\gamma$ term $0$ if $S_{t+1}$ is terminal). The loss is again $J_t(\mathbf{w})=\tfrac12\bigl(U_t-\mathbf{w}^\top\mathbf{x}_t\bigr)^2$. Full $\nabla_{\mathbf{w}} J_t$ keeps $\nabla_{\mathbf{w}} U_t=\gamma\mathbf{x}(S_{t+1},A_{t+1})$ when $S_{t+1}$ is nonterminal. Semi-gradient sets $\nabla_{\mathbf{w}} U_t=\mathbf{0}$ and yields

\begin{equation}
\mathbf{w}
\leftarrow
\mathbf{w}
+
\alpha
\Bigl[
R_{t+1}
+
\gamma\mathbf{w}^\top\mathbf{x}(S_{t+1},A_{t+1})
-
\mathbf{w}^\top\mathbf{x}_t
\Bigr]
\mathbf{x}_t.
\tag{7.4}
\end{equation}

Equivalently, $\mathbf{w}\leftarrow\mathbf{w}+\alpha\bigl[R_{t+1}+\gamma\hat{Q}(S_{t+1},A_{t+1};\mathbf{w})-\hat{Q}(S_t,A_t;\mathbf{w})\bigr]\mathbf{x}(S_t,A_t)$. Data and backup follow the same policy, so (7.4) is on-policy. As in the off-policy case, $U_t$ still \textbf{moves} every step after $\mathbf{w}$ updates. Section 7.2 discusses why (7.3) can diverge (deadly triad) while (7.4) is stable under standard conditions.

\begin{implnote}[Algorithm (semi-gradient SARSA with linear FA)]
Input: feature map $\mathbf{x}(s,a)$; step size $\alpha$; discount $\gamma$; exploration (e.g. $\varepsilon$-greedy).

Output: $\mathbf{w}$ such that $\hat{Q}(s,a;\mathbf{w})=\mathbf{w}^\top\mathbf{x}(s,a)\approx Q^\pi(s,a)$ for the behavior $\pi$; greedy readout $\pi'(s)=\arg\max_a\hat{Q}(s,a;\mathbf{w})$ optional.

\begin{enumerate}
\item Initialize $\mathbf{w}$ arbitrarily (e.g. $\mathbf{0}$).
\item Loop for each episode:
\begin{itemize}
\item Initialize $S$; choose $A$ from $S$ using a policy based on $\hat{Q}$ (e.g. $\varepsilon$-greedy).
\item For each step of the episode until $S$ is terminal:
\begin{itemize}
\item Take $A$; observe $R$, $S'$.
\item $\mathbf{x}\leftarrow\mathbf{x}(S,A)$.
\item If $S'$ is terminal:
\begin{itemize}
\item $U\leftarrow R$
\item $A'\leftarrow$ undefined (unused)
\end{itemize}
\item Else:
\begin{itemize}
\item Choose $A'$ from $S'$ using the same behavior policy.
\item $U\leftarrow R+\gamma\mathbf{w}^\top\mathbf{x}(S',A')$
\end{itemize}
\item $\mathbf{w}\leftarrow\mathbf{w}+\alpha\bigl(U-\mathbf{w}^\top\mathbf{x}\bigr)\,\mathbf{x}$.
\item $S\leftarrow S'$; $A\leftarrow A'$ (if $S'$ nonterminal).
\end{itemize}
\end{itemize}
\item Return $\mathbf{w}$ (and optional greedy readout of $\hat{Q}$).
\end{enumerate}

\end{implnote}

The updates above are \textbf{stochastic}. Each step uses one transition to form $U_t$ and the (full or semi) gradient of $J_t$, so no expectation appears inside the loss. A different objective is the squared \textbf{mean} residual $\tfrac12\bigl(\mathbb{E}[U\mid s,a]-\hat{Q}(s,a;\mathbf{w})\bigr)^2$ at a fixed pair. Its gradient is a product of two expectations and needs two independent draws of $(R,S')$---\textbf{double sampling}. That is not $\nabla_{\mathbf{w}} J_t$; online trajectories usually provide only one next state, so the algorithms use the stochastic $J_t$ updates.

\lstinputlisting{../code/ch07/linear_fa.py}


\section{The Deadly Triad}

The semi-gradient updates (7.3) and (7.4) share the same skeleton $\mathbf{w}\leftarrow\mathbf{w}+\alpha\,\delta_t\,\mathbf{x}_t$, with $\delta_t=U_t-\mathbf{w}^\top\mathbf{x}_t$ and a \textbf{shared} $\mathbf{w}$. Both use function approximation and bootstrapping. They differ in the third ingredient: (7.3) is off-policy; (7.4) is on-policy. That difference is the \textbf{deadly triad}.

\subsection{Three ingredients}

\begin{minipage}{\textwidth}
Instability can arise when these three are combined:
\begin{enumerate}
\item Function approximation (shared $\mathbf{w}$ couples many pairs)
\item Bootstrapping ($U_t$ depends on the current estimate)
\item Off-policy learning (data distribution $\neq$ backup's policy)
\end{enumerate}
\end{minipage}

Drop any one and the standard danger goes away. Tabular methods have no shared $\mathbf{w}$. Monte Carlo (no bootstrap) does not feed the estimate back into $U_t$. On-policy methods keep the distribution of updates matched to the policy in the backup.

\subsection{Off-policy divergence}

Off-policy semi-gradient FA---Q-learning FA (7.3) is the control example---can send $\|\mathbf{w}\|\to\infty$. The issue is a \textbf{mismatch}: behavior generates which pairs are updated, while the backup (e.g. $\max_{a'}$) pretends a different continuation. With a shared $\mathbf{w}$, correcting one pair moves predictions at others; the bootstrap can then amplify errors instead of damping them.

\citetxt{baird1995}{Baird, 1995} gave a standard example of that failure: an MDP with six states and reward $0$ on every transition, so the true values are all $0$, and a linear approximator that can represent those values. Shared features couple the states. Under direct semi-gradient updates, run equally often on every transition, the weights can still grow without bound. Baird also proposed residual-gradient methods that minimize the squared Bellman residual; this book keeps the semi-gradient updates (7.3)--(7.4) and does not develop those residual algorithms.

To see the shared-weight mechanism explicitly, reduce to two states (prediction form; $R=0$, $\gamma>1/2$). Set $\hat{V}(s_1)=w$ and $\hat{V}(s_2)=2w$, and update only at $s_1$ while the backup uses $s_2$:
\[
w
\leftarrow
w+\alpha\bigl(\gamma\cdot 2w-w\bigr)
=
\bigl(1+\alpha(2\gamma-1)\bigr)\,w.
\]
The factor is strictly larger than $1$, so $|w|\to\infty$. The same kind of mismatch sits inside (7.3) whenever exploring behavior disagrees with the greedy backup. One such instance is enough to show that off-policy linear semi-gradient is \textbf{not} generally safe.

\subsection{On-policy stability}

On-policy SARSA FA (7.4) keeps function approximation and bootstrapping, but \textbf{not} off-policy learning: $A_t$ and $A_{t+1}$ come from the same $\pi$, so the pairs that are updated and the pairs that appear in the bootstrap share one distribution.

Under standard conditions (linear features of full rank under $\pi$, bounded rewards, diminishing step sizes), on-policy linear TD and SARSA FA converge to a finite $\mathbf{w}$ (\citetxt{tsitsiklis1997}{Tsitsiklis and Van Roy, 1997}; see also \citetxt{suttonbarto2018}{Sutton and Barto}). The proof of that convergence is beyond the scope of this book. The takeaway here: (7.4) avoids the deadly triad's third vertex; (7.3) can hit all three.

\subsection{Summary}

\begin{center}
\begin{tabular}{lllll}
\toprule
Update & FA & Bootstrap & Off-policy & Deadly triad? \\
\midrule
(7.3) Q-learning FA & yes & yes & yes & can diverge \\
(7.4) SARSA FA & yes & yes & no & stable (linear, standard conditions) \\
\bottomrule
\end{tabular}
\end{center}

Chapter 7 used \textbf{linear} FA. Neural networks and other nonlinear approximators are taken up next: they still face the triad when paired with an off-policy bootstrap, and need practical stabilizers.
